%!latex
\documentclass[12pt, a4paper, leqno, twoside]{book}

\usepackage{%
  circuitikz,
  fancyhdr,
  geometry,
  graphicx,
  hyperref,
  newtxtext,
  newtxmath,
  pgfplots,
  setspace,
  titlesec,
  tikz,
  xcolor,
}

\usepackage[brazil]{babel}

% Sugerido por pelo pacote pgfplots
\pgfplotsset{compat=1.18}

\usetikzlibrary{math}

% Espaçamento de letra para o cabeçalho dos capítulos
\usepackage[letterspace=300]{microtype}

% Formatação do título dos capítulos
\titleformat{\chapter}[display]
{}%
{
   \centering\fontsize{80pt}{80pt}%
   {\selectfont\thechapter}
}%
{1em}% 
{\centering\scshape\lsstyle\MakeUppercase}%
[]

% Formatação do título das se\c c\~oes e subseções
\titleformat{\section}{\centering\large\scshape}{\thesection}{1em}{}[]
\titleformat{\subsection}{\centering\scshape}{\thesubsection}{1em}{}[]

% Estilo dos cabeçalhos das páginas
\renewcommand{\chaptermark}[1]{\markboth{#1}{}}
\renewcommand{\sectionmark}[1]{\markright{\thesection\ #1}}
\renewcommand{\headrulewidth}{0ex}
\fancypagestyle{est}{%
  \def\fontface{\bf}
  \fancyhf{}
  \fancyhead[CE]{\fontface\sc\leftmark}
  \fancyhead[CO]{\fontface\sc\rightmark}
  \fancyhead[LE,RO]{\fontface\thepage}
  \setlength{\headheight}{16pt}
  \addtolength{\topmargin}{-1pt}

}

% Numeração dos capítulos em algarismos romanos
\def\theHchapter{\arabic{chapter}\thechapter}
%\def\thechapter{\Roman{chapter}}

% Remover páginas em branco
\let\cleardoublepage=\clearpage

% Geometria da página
\geometry{top=3cm, bottom=2cm, left=3cm, right=2cm}

% Modificar o estilo da numeração das notas de rodapé
\renewcommand{\thefootnote}{\bfseries\arabic{footnote}}

% Definições gerais
\def\tpfc{\scshape}

\def\localData{%
  \parbox[c][0cm][c]{3.5cm}{
    \sc
    Jo\~ao Pessoa\hfill -\hfill PB\\
    \footnotesize
    \today
  }
}

\def\parteSuperior{
  \parbox[c][0cm][c]{10cm}{
    {\scshape\scriptsize%
      {INSTITUTO\hfill FEDERAL\hfill DE\hfill EDUCA\c C\~AO,\hfill CI\^ENCIA\hfill E\hfill TECNOLOGIA DA PARA\'IBA}\\
      {
        \fontsize{12}{12}\selectfont
        {\noindent CURSO\hfill T\'ECNICO\hfill SUBSEQUENTE\hfill EM\hfill ELETR\^ONICA}\\
        \large
        {\noindent TRABALHO\hfill DE\hfill CONCLUS\~AO\hfill DE\hfill CURSO}\\
      }
    }
  }
  \vfill 
  {\lsstyle\Large\scshape%
    PROJETO }
  \vskip .5cm
  \textit{por}
  \vskip .5cm
  {\tpfc Harllen Ara\'ujo de Sena}
  \vskip .15cm
  {\it e}
  \vskip .15cm
  {\tpfc Henrique Cirilo Costa}
  \vskip .5cm
  \textit{orientado pelo}
  \vskip .5cm
  {\tpfc Prof. Dr. C\'icero Alisson dos Santos}
  \vskip .5cm
}

\def\ohm{\,\Omega}
\def\ampop{amp-op}
\def\db{{\rm dB}}
\def\cmrr{{\rm CMRR}}

\def\[{\begin{equation}}
\def\]{\end{equation}}

\begin{document}

  \section{Entrada desbalanceada}
  \begin{figure}[!ht]
    \centering
    \begin{circuitikz}[line width=.7]
      \ctikzset{
        european,
      }
      \def\dmh{1.5}
      \def\dmv{3.6}
      \coordinate (K1) at (0,0);
      \coordinate (GND) at (2*\dmh,-\dmv);
      \draw (GND) to [short] (0,-\dmv) to [short,-o] (K1);
      \draw [thick, fill=white] (K1) coordinate (K1) circle [radius=1.5ex];
      \draw (K1) node [above=.2cm] {$K_1$} to [R, l=$R_1$, a={47R}, o-] (2*\dmh,0) to [C, l_={$C_1$}, a^={100p}, *-*] (GND);
      \draw (2*\dmh,0) to [short] ++(2*\dmh,0) coordinate (C2) to [R, l_=$R_2$, a^={220k}, *-] ++ (0,-\dmv) to [short] (2*\dmh,-\dmv);
      \draw (C2) to [short, -o] ++(0.7*\dmh,0) node [right] {DESBAL} ;
      \draw (GND) to ++ (0,-.3) node[tground] {};
    \end{circuitikz}
    \caption{Esquem\'atico da entrada desbalanceada.}
  \end{figure}
  Segue uma an\'alise AC da sa\'ida da entrada desbalanceada (de $1\,{\rm MHz}$ at\'e $100\,{\rm MHz}$):
  \begin{figure}[!ht]
    \centering
    \begin{tikzpicture}[scale=1]
      \begin{axis}[
        domain=10e6:10e8, 
        grid=both,
        grid style={black!20},
        scale only axis,
        xlabel={Frequ\^encia em Hz},
        xmax=10e7,
        xmin=10e5,
        xmode=log,
        ylabel={Amplitude em dB},
        ytick distance=1,
        width=12cm,
        height=7cm,
      ]
        \addplot [
          color=black, 
          smooth,
          ultra thick
        ] table {../amplificador_5532/entrada_desbalanceada.data};

        \addplot[
          domain=10e5:10e8, 
          samples=2,
        ] {-3};

        \node [
          pin={
            [fill=white, draw=black]
            265:{$f_c\approx34.6\,\rm MHz$}
          },
          circle,
          draw=black, 
          fill=white,
          inner sep=1.5pt
        ] at (axis cs:34e6,-3) {};

      \end{axis}
    \end{tikzpicture}
    \caption{An\'alise AC da sa\'ida da entrada desbalanceada.}
  \end{figure}

  \begin{figure}[!ht]
    \centering
    \begin{tikzpicture}[scale=1]
      \def\mw{1.7cm}
      \def\prop{1.2}
      \begin{axis}[
        domain=10e5:10e8, 
        grid=both,
        grid style={black!20},
        scale only axis,
        xlabel={Frequ\^encia em Hz},
        xmin=10e3,
        xmax=10e7,
        xmode=log,
        ylabel={Amplitude em dB},
        ytick distance=20,
        width=\prop*12cm,
        height=\prop*7cm,
      ]

        \addplot [
          thick,
          color=black, 
          smooth,
        ] table {../amplificador_5532/estagio_de_ganho.data};

        \addplot[
          domain=10e3:10e7, 
          samples=2,
        ] {-3} ;

        \node [
          anchor=east,
          draw=black,
          fill=white, 
          inner sep=2pt,
          minimum width=\mw,
        ] at (axis cs:10e7,-3) {$-3\,\rm dB$};


        \addplot[
          domain=10e3:10e7, 
          samples=2,
        ] {22.7} ;

        \node [
          anchor=east,
          fill=white, 
          draw=black,
          inner sep=2pt,
          minimum width=\mw,
        ] at (axis cs:10e7,22.7) {$22.7\,\rm dB$};



        %\node [
        %  pin={
        %    [fill=white, draw=black, rounded corners=2pt]
        %    225:{$f_c\approx34.6\,\rm MHz$}
        %  },
        %  circle,
        %  draw=black, 
        %  fill=red,
        %  inner sep=1.5pt
        %] at (axis cs:34e6,-3) {};

      \end{axis}
    \end{tikzpicture}
    \caption{An\'alise AC da sa\'ida do est\'agio de ganho.}
  \end{figure}
  \begin{figure}[]
    \centering
    \begin{circuitikz}[line width=.7, scale=.8, rotate=90, transform shape]
      \ctikzset{
        european,
      }
      \def\dmh{1.5}
      \def\dmv{3.6}
      \draw (0,0) node [left] {\it IN} to [C, l={$C_2$}, a={2u2}, o-] ++(1.5*\dmh,0) coordinate (C1) to [R, text width=1cm, align=center, l_={$R_3$ 47k},-] ++(0,-2*\dmv) to [short, -*] ++(\dmh/2,0) coordinate (GND);

      % IC1B
      \draw (C1) to [short] ++(2*\dmh,0) node [op amp, anchor=+, noinv input up] (IC1B) {IC1B};
      \draw (IC1B.-) to [short] ++ (-1.5*\dmh,0) to [short] ++(0,-\dmv/2) coordinate (C2)  to [short] ++(0,-\dmv/5) to [R, text width=1cm, align=center, l={$R_6$ 910R}] (GND) to [short] ++(0,-.5) node [tground] {};
      \draw (C2) to [short, *-] ++(\dmh,0) coordinate (C3) to [R, l^={$R_7$}, a={2k2}, -*] (C2 -| IC1B.out) to [short, -*] (IC1B.out) to [R, l={$R_9$}, a={10R}] ++(1.5*\dmh,0) coordinate (IC2ANIN); 
      \draw (C3) to [short, *-] ++ (0,-.6*\dmv) coordinate (C4) to [C, l={$C_4$}, a={33p}] (C4 -| IC1B.out) to [short] (C2 -| IC1B.out);

      % IC2B
      \draw (C1) to [short, *-] ++(0,2*\dmv) to [short] ++(2*\dmh,0) node [op amp, anchor=+, noinv input up] (IC2B) {IC2B};
      \draw (IC2B.-) to [short] ++ (-1.5*\dmh,0) to [short] ++(0,-\dmv/2) coordinate (C5)  to [short, *-] ++(0,-\dmv/5) to [R, l^={$R_4$\\ 910R}, text width=1cm, align=center] ++(0,-.7*\dmv) to [short] ++(0,-.5) node [tground] {};
      \draw (C5) to [short] ++(\dmh,0) coordinate (C6) to [R, l^={$R_5$}, a={2k2}, *-*] (C5 -| IC2B.out) to [short] (IC2B.out) to [R, l={$R_8$}, a={10R}, *-] ++(1.5*\dmh,0) to [short, -*] (IC2ANIN); 
      \draw (C6) to [short] ++ (0,-.6*\dmv) coordinate (C7) to [C, l={$C_3$}, a={33p}] (C7 -| IC1B.out) to [short] (C5 -| IC1B.out);

      % IC2A
      \draw (IC2ANIN) to [short] ++(1.5*\dmh,0) node [op amp, anchor=+, noinv input up] (IC2A) {IC2A};
      \draw (IC2A.-) to [short] ++ (-1.5*\dmh,0) to [short] ++(0,-\dmv/2) coordinate (C8)  to [short] ++(0,-\dmv/5) to [R,l_={$R_{10}$}, a^={2k2}] (GND -| C8) to [short] ++(0,-.5) node [tground] {};
      \draw (C8) to [short, *-] ++(\dmh,0) coordinate (C9) to [R, l^={$R_{11}$}, a={2k2}, *-*] (C8 -| IC2A.out) to [short, -*] (IC2A.out); 
      \draw (C9) to [short] ++ (0,-.6*\dmv) coordinate (C10) to [C, l={$C_6$}, a={33p}] (C10 -| IC2A.out) to [short] (C9 -| IC2A.out);

      % IC2B
      \draw (IC2A.out) to [short] ++ (\dmh/2,0) node [op amp, anchor=+, noinv input up] (IC2B) {IC2B};
      \draw (IC2B.-) to [short] ++(0,-\dmv/3) coordinate (C11) to [short] (C11 -| IC2B.out) to [short] (IC2B.out) to [R, l={$R_{12}$}, a={1k}, *-*] ++ (1.5*\dmh,0) coordinate (IC3BINV);

      % IC3B
      \draw (IC3BINV) node [op amp, anchor=-] (IC3B) {IC3B};
      \draw (IC3B.+) to [short] ++ (-\dmh/10,0) to [short] ++ (0,-.5) node [tground] {};
      \draw (IC3B.out) to [short] ++ (0,.6*\dmv) coordinate (C12) to [C, l_={$C_8$}, a^={33p}, -*] (C12 -| IC3BINV);
      \draw (IC3B.out) to [R, l={$R_{14}$}, a={47R}, *-*] ++(1.5*\dmh,0) coordinate (C14) to [short] ++(0,\dmv) coordinate (C13) to [R, l_={$R_{13}$}, a^={2k}, *-] (C13 -| IC3BINV) to [short] (IC3BINV);
      \draw (C14) to [short, -o] ++(\dmh/2,0) node [text width=2cm, align=center, right] {$K3$\\\it BRIDGE};
      \draw (C14) to [short] ++ (0, \dmv);

      % IC3A
      \draw (IC3BINV) ++ (0,4.5*\dmh) coordinate (IC3AINV) node [op amp, anchor=-] (IC3A) {IC3A};
      \draw (IC3A.+) to [short] ++ (-\dmh/10,0) to [short] ++ (0,-.5) node [tground] {};
      \draw (IC3A.out) to [short] ++ (0,.6*\dmv) coordinate (C15) to [C, l_={$C_9$}, a^={33p}, -*] (C15 -| IC3AINV);
      \draw (IC3A.out) to [R, l={$R_{17}$}, a={47R}, *-*] ++(1.5*\dmh,0) coordinate (C16) to [short] ++(0,\dmv) coordinate (C15) to [R, l_={$R_{16}$}, a^={2k2}] (C15 -| IC3AINV) to [short] (IC3AINV);
      \draw (C13) to [short] ++(0,\dmv/3) to [short] ++(-5*\dmh,0) coordinate (C17) to [short] (C17 |- IC3AINV) to [R, l^={$R_{15}$}, a_={2k2}, -*] (IC3AINV);
      \draw (C16) to [short, -o] ++(\dmh/2,0) node [text width=2cm, align=center, right] {$K4$\\\it PARALLEL}; 


      %\coordinate (K1) at (0,0);
      %\coordinate (GND) at (2*\dmh,-\dmv);
      %\draw (GND) to [short] (0,-\dmv) to [short,-o] (K1);
      %\draw [thick, fill=white] (K1) coordinate (K1) circle [radius=1.5ex];
      %\draw (K1) node [above=.2cm] {$K_1$} to [R=$R_1$,o-] (2*\dmh,0) to [C={$\displaystyle C_1\atop{\rm 100p}$}, *-*] (GND);
      %\draw (2*\dmh,0) to [short] ++(2*\dmh,0) coordinate (C2) node [above] {DESBAL} to [R=$R_2$, *-] ++ (0,-\dmv) to [short] (2*\dmh,-\dmv);
      %\draw (GND) to ++ (0,-.5) node[tground] {};
    \end{circuitikz}
    \caption{Esquem\'atico do est\'agio de ganho.}
  \end{figure}

  \begin{figure}[!ht]
    \begin{circuitikz}[scale=.9, transform shape]
      \ctikzset{
        european,
      }
      \def\dmh{1.5}
      \def\dmv{3.6}
     \foreach \y in {0,...,7} {
       \draw (0,\y*.8*\dmv) to [short] ++(\dmh,0) node [op amp, anchor=+, noinv input up] (IC\y) {IC\y};
       \draw (IC\y.-) to [short] ++(0,-.3*\dmv) coordinate (C1) to [short] (C1 -| IC\y.out) to [short, -*] (IC\y.out) to [R, l={R}, a={1R}] ++(1.5*\dmh,0) to [short] ++ (0,-.5*\dmv);
     }
    \end{circuitikz}
  \end{figure}

\end{document}
